\documentclass[11pt]{article}
% Shared preamble for private TrustAndReview manuscript drafts.
\usepackage[margin=1in]{geometry}
\usepackage[T1]{fontenc}
\usepackage[utf8]{inputenc}
\usepackage{lmodern}
\usepackage{microtype}
\usepackage{amsmath,amssymb,mathtools}
\usepackage{booktabs,tabularx,longtable}
\usepackage{enumitem}
\usepackage{xcolor}
\usepackage{hyperref}
\usepackage[round,authoryear]{natbib}
\hypersetup{colorlinks=true,linkcolor=blue!50!black,citecolor=blue!50!black,urlcolor=blue!60!black}
\setlist{nosep}

\newcommand{\draftstatus}[1]{%
  \begin{center}\fcolorbox{black}{yellow!15}{\parbox{0.92\linewidth}{\textbf{Private draft status:} #1}}\end{center}}
\newenvironment{evidencegate}[1][Evidence required before promotion]
  {\begin{quote}\color{red!60!black}\textbf{#1.} }
  {\end{quote}}
\newcommand{\designchoice}[1]{\textbf{Design choice:} #1}
\newcommand{\openhypothesis}[1]{\textbf{Open hypothesis:} #1}
\newcommand{\implementationobservation}[1]{\textbf{Implementation observation:} #1}


\usepackage{array}
\usepackage{caption}
\usepackage{float}
\usepackage{graphicx}
\usepackage{url}
\newcommand{\tdrg}{TDRG}
\newcommand{\mcrp}{MCRP}
\newcommand{\testid}[1]{\texttt{#1}}
\newcommand{\roadmap}{\textsf{Roadmap}}
\newcommand{\verifiedslice}{\textsf{Verified v0.1 slice}}

\title{Plural Trust Without a Global Reviewer Score:\\
A Conformance Profile and Adversarial Evaluation Plan for Machine-Verifiable Scientific Review}
\author{Private working draft---authors and contribution record pending}
\date{24 August 2026}

\begin{document}
\maketitle

\draftstatus{Protocol/evaluation first draft. The only implementation evidence admitted here is local execution of 11 tests at tooling revision \texttt{9da0cd4}. TM-01--TM-18 are otherwise a roadmap. This draft is not approved for submission or public release.}

\begin{abstract}
Machine-verifiable scientific records can preserve claims, evidence, workflows,
and review events, but they do not answer a separate governance question: whose
review should a particular community solicit or rely upon? A single reviewer
reputation score collapses topic, capability, policy, time, and relying-party
context into a common failure domain. We specify the Trust-Domain and
Review-Governance profile (\tdrg{}), a draft composition of prior mechanisms that
keeps a shared scientific and review event record separate from community-relative
routing views. The profile defines typed events, root- and topic-scoped views,
importer-approved bridges, forkable policy overlays, explicit authority boundaries,
and evidence-promotion rules. We also define 18 stable conformance and adversarial
test families and a frozen comparison among five routing arms. A deliberately
small Python reference artifact currently demonstrates only local-view evaluation,
event-type exclusion, exact topic and bridge checks, time filtering, root relativity,
two malformed-input rejections, and a multiaxis reproduction-resource classifier;
11 unit tests passed in a recorded local rerun. The artifact does not implement
signatures, uniqueness, Sybil or collusion defenses, case-persona anonymity,
appeals, fork reconciliation, standards round trips, or comparative attack
experiments. Conformance is therefore a statement about named behavior in named
fixtures, not evidence of reviewer competence, anonymity, scientific correctness,
or social benefit.
\end{abstract}

\section{Introduction}

Scientific reproducibility and scientific review are related but non-equivalent.
A reproducibility record can bind a claim to data, code, configuration, environment,
compute, and validation evidence. It can record who or what executed a workflow and
which review object challenged a claim. None of those facts, by itself, determines
which future reviewer is competent, independent, available, or trusted by a given
community. Conversely, a reviewer-routing system cannot establish that a scientific
claim is true. Treating either record as a substitute for the other turns
administrative standing into scientific authority.

Existing work already supplies many of the relevant parts. Repository-centered
peer review and graph-based reviewer discovery predate the present proposal
\citep{rodriguez2006convergence}. Fair and randomized reviewer assignment have
formal and empirical foundations \citep{stelmakh2021peerreview4all,jecmen2020mitigating}.
Global recursive reputation, root-relative attack-resistant trust, and local trust
metrics are established ideas \citep{kamvar2003eigentrust,levien1998attack,massa2005controversial}.
Anonymous credentials and anonymous reputation have also been applied to review or
reputation settings \citep{naessens2006fair,zhai2016anonrep}. Decentralized open-review
architectures explicitly address interoperability and governance
\citep{tenoriofornes2021decentralizing}. The contribution claimed here is therefore
not invention of these mechanisms.

The narrower problem is compositional. A machine-verifiable scientific record needs
to remain usable when communities disagree about roots, competency criteria,
conflicts, privacy policy, or sanctions. It also needs to prevent an evaluation of a
review from silently becoming a scientific disposition, a conduct finding, or a
portable score on the reviewer. A protocol must expose the trust assumptions on
which a local routing view depends, while allowing communities to fork their policy
without rewriting the underlying scientific and review history.

This paper makes three bounded contributions. First, it defines \tdrg{} as a
governance and conformance profile over a shared event record and plural policy
views. Second, it supplies a threat-driven test contract with stable identifiers
\testid{TM-01}--\testid{TM-18}, including a preregisterable A0--A4 comparative
design. Third, it records the exact behavior of a small reference implementation
and separates that evidence from the unimplemented roadmap. The current evidence
is intentionally modest: 11 local unit tests on a synthetic fixture. The design is
useful only if its future claims can fail visibly.

The Minimum Credible Reproducibility Protocol (MCRP) core governs the claim/evidence
scientific record. \tdrg{} is optional and separately versioned; it remains outside
MCRP core conformance until its privacy, accountability, interoperability, replay,
and empirical promotion gates pass. Paper~01 maps the literature and cross-layer
interfaces, Paper~02 owns the formal model and A0--A4 research program, and this
paper owns the proposed normative exchange objects and conformance contract.

\designchoice{The protocol has no universal actor score and no global
\texttt{trusted=true}. A view is an artifact produced for a relying party under a
declared domain version, topic, capability, role, policy, evidence snapshot, and
time.}

The remainder of the paper first fixes terminology and the novelty boundary,
then states the threat model and protocol objects. It next distinguishes the
verified v0.1 implementation slice from the test roadmap, defines conformance and
comparative evaluation, addresses privacy and standards interoperability, and
closes with reproducibility obligations and stop rules.

\section{Positioning and novelty boundary}

\subsection{Routing, reputation, and review infrastructure are prior art}

The phrase ``decentralized peer review'' covers several different functions:
transporting review requests, selecting reviewers, publishing review objects,
aggregating feedback, rewarding participants, and governing an editorial service.
\citet{rodriguez2006convergence} described a repository-centric model with review
metadata and social-network-based reviewer selection. \citet{tenoriofornes2021decentralizing}
later developed an interoperable decentralized open-review architecture and
evaluated prototypes and operating costs. These works preclude a broad novelty
claim based merely on moving review outside a journal or using a graph.

Assignment is also its own mature problem. PeerReview4All optimizes the quality of
the worst-served paper and studies statistical accuracy under explicit models
\citep{stelmakh2021peerreview4all}. Randomized assignment can constrain the
probability of specific reviewer--paper matches and address manipulation and
assignment-based deanonymization under its model \citep{jecmen2020mitigating}.
\tdrg{} should consume an appropriate assignment method after eligibility is
computed; it should not rebrand routing randomization as a trust innovation.

The trust literature gives equally strong warnings against novelty inflation.
EigenTrust produces a global recursive score using pre-trusted peers
\citep{kamvar2003eigentrust}. \citet{levien1998attack} analyze root-relative
certificate-graph trust and a capacity-constrained max-flow construction under
explicit node- and edge-attack models. \citet{massa2005controversial} compare local
and global trust for controversial users in a recommender setting. These results
motivate comparison arms; they do not prove that the same mechanisms improve
scientific review.

Privacy mechanisms are likewise established. \citet{naessens2006fair} combine
anonymous credentials, topic-specific reputation, and accountability in a submission
and review system. AnonRep demonstrates anonymous reputation under an explicit
multi-provider trust model \citep{zhai2016anonrep}. Consequently, a fresh case
persona is a protocol choice to be threat-tested, not a claim that anonymous
reputation is new or solved.

\subsection{Residual contribution}

The proposed residual contribution is a constrained synthesis:

\begin{enumerate}
  \item scientific and review events remain shared, immutable inputs rather than
  being duplicated inside a reputation system;
  \item routing views are relative to declared roots, topics, capabilities,
  policies, evidence snapshots, and times;
  \item review evaluation, future reliance, scientific disposition, allegation,
  and conduct adjudication are distinct event types;
  \item cross-domain authority requires an explicit, attenuated bridge accepted by
  the importing domain; and
  \item governance forks create new overlays without deleting common scientific or
  review objects.
\end{enumerate}

This contribution survives only as far as the profile is precise, interoperable,
and empirically distinguishable from simpler baselines. If A4 does not provide a
useful frontier over A0--A3, \tdrg{} should be reported as an interoperability
vocabulary and governance checklist, not promoted as a superior review mechanism.

\section{Terminology and evidence boundary}

We distinguish five properties that are often compressed into ``reproducible.''
\emph{Rerenderability} rebuilds presentation artifacts from supplied derived data.
\emph{Computational repeatability} repeats a declared computation with fixed inputs
and a sufficiently equivalent environment. \emph{Computational reproducibility}
recovers claim-relevant findings while independently reconstructing a meaningful
portion of the computational chain. \emph{Independent replication} obtains new
evidence through an independently implemented or newly observed route.
\emph{Scientific validity} is the substantive judgment that the inference is
adequate for its claim. A Makefile can contribute to the first two without reaching
the latter three.

The same discipline applies to review governance. A \emph{review event} records
an evaluation of a scientific object. A \emph{reliance event} expresses willingness
to route work to, supervise, or otherwise rely on an actor under a scope. A
\emph{scientific disposition} records what a decision process did with a claim. An
\emph{adverse event} records an allegation or routing precaution. A
\emph{conduct finding} is the result of an authorized process. These events may
refer to one another, but one must not be silently manufactured from another.

We use three evidence labels throughout:

\begin{description}
  \item[\verifiedslice] behavior exercised by a named test in tooling revision
  \texttt{9da0cd4\ldots} and rerun locally on 24 August 2026;
  \item[\roadmap] a specified object, test, or comparison without admissible result;
  \item[open hypothesis] a falsifiable comparative proposition awaiting frozen
  inputs, execution, uncertainty analysis, and adverse-result reporting.
\end{description}

\implementationobservation{The v0.1 tooling repository contains a Python
standard-library evaluator, one documentary JSON Schema, one synthetic
gravitational-wave-inspired view fixture, a resource-profile classifier, and 11
unit tests. All 11 tests passed in the local rerun recorded for this draft.}

This observation supports only the precise behaviors listed in
Section~\ref{sec:implementation}. It does not promote an entire test-matrix row
merely because one subcase is present.

\section{Requirements and threat model}
\label{sec:threat}

\subsection{Adversaries and failure domains}

The protocol must be interpretable under malicious, negligent, poorly calibrated,
coerced, and conflicted behavior. The relevant adversary is not limited to a
pseudonymous outsider. It may include a credentialed review ring, a captured trust
root, a line manager pressuring reviewers, a registrar colluding with a publisher,
an assignment service manipulating exposure, or an evaluator accepting malicious
input. Staleness and operational failure are also adversarial in effect even when
no actor is malicious.

Table~\ref{tab:threats} maps representative threats to the property under test and
the residual claim boundary. ``Trusted operator'' means an actor whose correct or
non-colluding behavior is assumed for that property; naming it exposes, rather than
eliminates, the trust boundary.

\begin{table}[H]
\centering
\caption{Threat-to-test map. Each row names a residual risk; no row is a claim that
the current implementation supplies the control.}
\label{tab:threats}
\small
\begin{tabularx}{\textwidth}{>{\raggedright\arraybackslash}p{0.19\textwidth} X p{0.16\textwidth} p{0.13\textwidth}}
\toprule
Threat & Protected property and assumed boundary & Test family & Residual risk \\
\midrule
Sybils and real-person collusion & Scoped influence; honest credential and control metadata & TM-07, TM-08, TM-18 & Colluding real people can pass uniqueness checks \\
Bad-mouthing and review rings & Typed, nontransitive adverse information; qualified dissent remains visible & TM-04, TM-05, TM-18 & Biased ground truth or captured adjudication \\
Cross-topic laundering & Importer consent, exact scope, attenuation, expiry & TM-06 & Ontology mismatch and legitimate multidisciplinary exclusion \\
Captured roots or boards & View binds roots; policy can fork without rewriting records & TM-07, TM-14 & A captured relying party may still choose bad roots \\
Management pressure & Separation of assignment, opening, conduct, and line authority & TM-12, TM-13 & Informal coercion outside recorded channels \\
Deanonymization & Declared adversary, measured linkage, claim downgrade & TM-10--TM-12 & Prose, timing, rare attributes, and small pools \\
Rollback and stale state & Explicit clocks, expiry, version binding, adverse freshness status & TM-15 & Partitions or compromised clock and evaluator sources \\
Malicious documents & Data cannot acquire policy or execution authority & TM-17 & Parser, model, and tool-chain vulnerabilities \\
\bottomrule
\end{tabularx}
\end{table}

\subsection{Security anti-claims}

A signature can bind bytes to a key under a chosen verification procedure. It does
not establish that the key belongs to one person, that the person is competent or
independent, that the signed statement is honest, or that the underlying science is
correct. Similarly, a graph path records a relation under a policy; it does not make
the endpoint a valid reviewer for every case. These distinctions are conformance
requirements because otherwise a technically valid artifact can carry an
epistemically invalid implication.

\begin{evidencegate}[Security and privacy promotion gate]
No signature, uniqueness, Sybil-resistance, case-persona unlinkability, opening, or
collusion-resistance result is admissible until a versioned implementation, declared
operator coalition, attack fixture, and report are available. The current artifact
contains none of these mechanisms.
\end{evidencegate}

\section{The \tdrg{} protocol profile}
\label{sec:protocol}

\subsection{Shared records and local views}

The profile separates two planes. The \emph{record plane} contains scientific
claims, evidence, execution provenance, reviews, responses, challenges,
supersessions, and governance events. The \emph{view plane} contains derived,
recomputable routing or supervision views. A view refers to immutable record-plane
objects and a frozen policy snapshot; it must not mutate its inputs.

Formally, a local view is represented as
\begin{equation}
V(S,d,z,c,r,q,P,E_{\leq\tau},t_{\mathrm{eval}})
 = F(S,d,z,c,r,q,P,E_{\leq\tau},t_{\mathrm{eval}}),
\end{equation}
where $S$ is the relying root set, $d$ is the domain-manifest version, $z$ is the
topic scope, $c$ is capability, $r$ is review role, $q$ is the case or reliance
question, $P$ is evaluator policy, $E_{\leq\tau}$ is the frozen typed-event cut,
and $t_{\mathrm{eval}}$ is the evaluation clock with
$\tau\leq t_{\mathrm{eval}}$. The distinct cut and clock permit historical replay
under explicitly selected historical or current freshness and expiry rules. This
equation is an interface, not a claim that one aggregation function is socially
optimal.

\designchoice{A conforming view reports every input binding and the policy-defined
reason that an event was included or excluded. Absence of a path remains
``unknown'' and must not become adverse evidence.}

\subsection{Versioned domain manifest}

A domain manifest identifies governance parentage, included and excluded topics,
capabilities, roles, recognized issuers, roots, evaluator and parameters, bridge
policy, staleness rules, assignment and newcomer policy, authority separation,
privacy claims, appeals, compensation, and fork/exit conditions. Material policy
change creates a new version. A mutable landing page is not sufficient because a
later observer must be able to recover the exact policy under which a view was
computed.

The manifest does not prove that its controllers are legitimate. It makes the
controller and claimed boundary inspectable. Independent actors may reject the
manifest, recognize only part of it, or import nothing.

\subsection{Identity scopes and authority separation}

The draft distinguishes an identity record, a domain account, and a review-case
persona. The public persona is intended to be fresh for the tuple consisting of
domain version, scientific-record digest, request, round, role, and policy version.
Private services may need to detect duplicate participation or process a qualified
appeal, but routine public linkage across cases is prohibited by design.

This division creates an explicit operator problem. Registrar, eligibility service,
assignment authority, publisher, reputation custodian, policy evaluator, conduct
adjudicator, and exceptional-opening trustees are separate roles. Combining roles
does not automatically make a deployment nonconforming, but the manifest must
state the combination and downgrade any privacy or independence claim that the
combination defeats.

\begin{evidencegate}[Persona gate]
The current implementation has no identity, credential, signature, case-persona,
duplicate-prevention, opening, or appeal code. This subsection specifies desired
semantics only; it makes no anonymity claim.
\end{evidencegate}

\subsection{Typed events and scientific disposition}

At minimum, the event vocabulary distinguishes:

\begin{itemize}
  \item a review object evaluating a claim, evidence item, workflow, or release;
  \item an evaluation of that review object;
  \item a scoped reliance delegation for future routing or supervision;
  \item a scientific disposition such as accepted, challenged, or unresolved;
  \item an allegation or precautionary routing event;
  \item an authorized conduct finding and any later correction; and
  \item a challenge, appeal, or supersession event.
\end{itemize}

Policies may read relationships among these types, but a conversion must be an
explicit, reviewable rule with an issuing authority and appeal path. In particular,
disagreement with a scientific conclusion cannot automatically produce a conduct
finding or an identity-opening request.

\subsection{Root-relative views}

The v0.1 prototype uses a deliberately simple reference function: among eligible
explicit positive trust events it computes the maximum product of edge weight and
exponential time-decay along paths up to a declared depth. For a path
$p=(e_1,\ldots,e_k)$ evaluated at time $t$,
\begin{equation}
s(p,t)=\prod_{e\in p} w_e\,2^{-\Delta t_e/h}, \qquad
S_R(x,t)=\max_{p:R\leadsto x}s(p,t),
\end{equation}
where $h$ is the declared half-life. The score is local to the root, topic, allowed
bridges, clock, and input events. Maximum-path semantics were chosen for
inspectability, not because evidence shows they are the best governance rule.

This prototype rule is neither EigenTrust nor the capacity-constrained max-flow
metric of \citet{levien1998attack}. Those mechanisms belong in comparison arms.
The prototype also has no negative-edge propagation because negative standing is
not implemented at all.

\subsection{Bridges and nested domains}

A bridge is an import authorization, not a declaration of global equivalence. It
names the exporting domain, importing domain, topic and capability scope,
attenuation or maximum authority, supervision requirement, issue and expiry time,
and policy version. The importer must approve it. A parent domain's recognition of
a child does not automatically export standing upward, downward, or sideways.
An ordinary bridge may map credential, capability, review-event, and evidence
vocabularies and may reference an MCRP claim identifier, but it cannot rewrite or
promote that claim. Scientific dispositions, sanctions, and negative standing do
not transfer by default; each requires a separately authorized typed-event mapping.

Nested domains have at most one normative governance parent but may recognize many
peer domains or competency imports. This separates constitutional inheritance from
multidisciplinary evidence. An exact topic match in the prototype is only a minimal
mechanism demonstration; production interoperability will require versioned
ontologies and explicit loss handling.

\subsection{Forks and fission}

A fork creates new roots, policy, and derived views while retaining identifiers for
the shared scientific and public review record. Private identity mappings, private
standing, and restricted evidence are not copied without an independent access
basis and consent. Reconciliation is represented by later recognition or bridge
events, not by erasing the divergence.

\designchoice{Forkability is a declared dependency-separation mechanism, not a
demonstrated containment of institutional capture or a guarantee that the
better-governed branch wins. It separates evaluator inputs only to the extent that
record identity, export, and independent operation remain available.}

\section{Reference implementation and exact evidence}
\label{sec:implementation}

\subsection{Artifact boundary}

The inspected prototype is private repository \texttt{trust-and-review-tooling},
revision \texttt{9da0cd4} (full digest recorded in the claim ledger). It uses the Python
standard library and contains three relevant components:

\begin{enumerate}
  \item a manual manifest validator and deterministic local-view evaluator;
  \item a documentary Draft 2020-12 JSON Schema for the view input; and
  \item a resource-profile validator/classifier with R0--R3 technical-access
  burden labels and four downselect modes.
\end{enumerate}

The JSON Schema is documentary in v0.1; no independent schema-engine validation
was run. The local command
\begin{quote}\ttfamily\small
PYTHONDONTWRITEBYTECODE=1 python3 -m unittest discover -s tests -v
\end{quote}
reported 11 tests run and 11 passing on 24 August 2026. The run was fast and uses
synthetic inputs. It is a local verification observation, not an immutable archival
release or independent reproduction.

\subsection{Subtest-level coverage}

Table~\ref{tab:implemented} is the complete admitted implementation result. A
partial mapping to a TM row does not mean the row passed.

\begin{longtable}{p{0.23\textwidth}p{0.17\textwidth}p{0.49\textwidth}}
\caption{Exact v0.1 test evidence. All observations are scoped to the named unit
test and revision.}\label{tab:implemented}\\
\toprule
Test & Matrix relation & Verified observation \\
\midrule
\endfirsthead
\toprule Test & Matrix relation & Verified observation \\
\midrule
\endhead
\path{test_expected_local_view_and_audit_reasons} & TM-04/06/15 partial & One fixture excludes a review-rating event as non-trust, denies one unapproved bridge, rejects one wrong-topic edge, and excludes one expired and one revoked edge; the expected actors and score ordering are present. \\
\path{test_bridge_is_importer_approved_and_topic_specific} & TM-06 partial & Removing the allowed bridge removes the bridged specialist from one view. \\
\path{test_view_is_root_relative} & TM-07 partial & Changing the root changes reachability in one fixture. No frontier-capacity or capture bound is tested. \\
\path{test_view_is_topic_colored} & TM-06 partial & Changing the exact topic changes included actors in one fixture. \\
\path{test_future_revocation_does_not_apply} & TM-15 partial & A revocation issued after the view clock does not remove its target in one fixture. \\
\path{test_output_is_deterministic_under_event_reordering} & TM-02 partial & Reversing event input order yields equal Python result objects in one process and fixture. This is not a cross-platform byte-identity result. \\
\path{test_malformed_and_negative_weight_are_rejected} & TM-01/17 partial & One negative weight raises \texttt{ProtocolError}. No document or agent prompt-injection path is tested. \\
\path{test_duplicate_event_ids_are_rejected} & TM-01/17 partial & One duplicate event identifier raises \texttt{ProtocolError}. \\
\path{test_hpc_profile_preserves_axes_and_labels_burden} & Reproducibility profile & One synthetic profile is classified R2; human-expert hours remain visible; a downselect is not labeled full reproduction. \\
\path{test_nonfull_mode_must_declare_excluded_claims} & Result promotion & One non-full mode with an empty excluded-claims list raises \texttt{ProtocolError}. \\
\path{test_portable_profile_is_level_zero} & Reproducibility profile & One modified lightweight profile is classified R0. \\
\bottomrule
\end{longtable}

\subsection{What the result does not support}

The current artifact does not supply or test: cryptographic signatures or
attestations; personhood or uniqueness; Sybil or collusion resistance; negative
edge containment; capacity bounds; control diversity; fair assignment; workload or
newcomer policy; case personas; anonymity or confidentiality measurement;
authority-collusion constraints; appeals or opening; forks or reconciliation;
standards projections; stale-evaluator rollback; sandboxing of malicious documents;
or A0--A4 experiments. It also does not validate a scientific claim or demonstrate
field efficacy. These omissions are not failed tests; they are absent mechanisms.

\begin{evidencegate}[Archival result gate]
Before the local observation becomes an archival implementation claim, create a
tagged release that binds protocol and schema digests, code revision, environment,
fixture set, clock, command, platform, test output, and artifact digest. Independent
rerun and license review remain required.
\end{evidencegate}

\section{Conformance and adversarial test contract}
\label{sec:tests}

Conformance is a tuple, not a badge:
\begin{equation}
C=(v_{protocol},v_{schema},v_{code},E,P,t,H,T,O),
\end{equation}
where $E$ is the frozen event set, $P$ the policy, $t$ the clock, $H$ the execution
environment and hardware class, $T$ the named tests, and $O$ their outputs. A claim
must point to the exact tuple and list exclusions.

Table~\ref{tab:matrix} defines stable test-family semantics. Rows may be split into
subtests but not silently redefined. Except for the explicit partial mappings in
Table~\ref{tab:implemented}, all entries are roadmap items.

\begin{longtable}{p{0.10\textwidth}p{0.27\textwidth}p{0.53\textwidth}}
\caption{TDRG conformance and adversarial roadmap.}\label{tab:matrix}\\
\toprule ID & Property & Minimum passing evidence \\
\midrule
\endfirsthead
\toprule ID & Property & Minimum passing evidence \\
\midrule
\endhead
TM-01 & Object/schema conformance & Valid objects accepted and minimally mutated invalid objects rejected with stable reasons for every normative type. \\
TM-02 & Deterministic replay & Frozen inputs repeated in clean environments produce canonical byte-identical output or a documented semantic digest equivalence. \\
TM-03 & Policy sensitivity & Changing exactly one declared policy input produces an attributable and explainable delta. \\
TM-04 & Event-type separation & Votes, review evaluations, dispositions, allegations, and findings cannot undergo forbidden automatic conversion. \\
TM-05 & Negative-edge containment & Adverse events do not propagate transitively or create enemy-of-enemy positive trust. \\
TM-06 & Topic/capability isolation & Transfer occurs only through accepted, unexpired, attenuated bridges with declared supervision. \\
TM-07 & Root relativity/frontier bounds & Outputs bind roots and malicious influence respects the declared capacity rule. \\
TM-08 & Control diversity & Employer, supervisor, funding, and infrastructure control are reported separately from graph-path diversity. \\
TM-09 & Fair routing/newcomers & Workload and newcomer policy are enforced without deterministic top-rank selection. \\
TM-10 & Case persona/replay & Cases use distinct public personas, preserve allowed within-case continuity, and reject duplicates. \\
TM-11 & Anonymity downgrade & Declared linkage attacks are measured and the claim downgrades when its threshold is missed. \\
TM-12 & Authority collusion & Capability and privacy loss under named operator coalitions match the manifest. \\
TM-13 & Appeals/opening & Appeals preserve the required privacy, corrections append, and unauthorized opening is rejected. \\
TM-14 & Fork/bridge governance & Shared identifiers survive policy divergence; private state is not copied without authority. \\
TM-15 & Staleness/rollback & Expiry, revocation delay, partitions, clocks, and old evaluators produce distinct, declared statuses. \\
TM-16 & Standards projection & Export/reimport preserves identity and claim scope, challenges, supersession, and policy meaning. \\
TM-17 & Agent-input safety & Untrusted content remains data and cannot change policy or execution authority in the declared environment. \\
TM-18 & A0--A4 comparison & Matched synthetic scenarios produce a complete bundle including uncertainty, subgroups, cost, and adverse outcomes. \\
\bottomrule
\end{longtable}

Passing TM-01--TM-17 would establish only tested implementation behavior for a
release and fixture set. TM-18 could support synthetic comparative claims. Neither
establishes benefit in an independently governed community or validity of any
reviewed science.

\section{Comparative evaluation}
\label{sec:evaluation}

\subsection{Frozen arms}

The comparison must hold actors, cases, planted defects, harmless deviations, false
alarms, unresolved disputes, event order, and adjudications fixed across arms. Once
an arm determines eligibility, the same fair or randomized assignment procedure
should be used so that assignment quality is not misattributed to the trust model.

\begin{table}[H]
\centering
\caption{A0--A4 routing arms. These are preregistration definitions, not implemented
comparisons.}
\label{tab:arms}
\small
\begin{tabularx}{\textwidth}{p{0.08\textwidth}p{0.32\textwidth}X}
\toprule Arm & Mechanism & Interpretive purpose \\
\midrule
A0 & Eligibility plus uniform random routing & Non-reputation baseline \\
A1 & Raw feedback average or vote count & Deliberately weak popularity baseline \\
A2 & Global EigenTrust/PageRank-like recursive score & Global recursive baseline \\
A3 & Topic-scoped personalized restart walk & Root-relative local baseline \\
A4 & A3 with frontier-capacity bounds, then control-diversity, workload, and newcomer constraints & Full proposed routing composition \\
\bottomrule
\end{tabularx}
\end{table}

Implementations of A2--A4 must be published precisely enough that an observed
difference is not an artifact of unequal tuning. A2 should be faithful to a named
global recursive method; A4's capacity component should state how it relates to,
and differs from, established max-flow trust models
\citep{kamvar2003eigentrust,levien1998attack}. The assignment layer should be
separately identified \citep{stelmakh2021peerreview4all,jecmen2020mitigating}.

\subsection{Scenarios and outcome measures}

The synthetic population should include overlapping topic layers, a
multidisciplinary case, sparse specialties, newcomers, narrow experts, calibrated
reviewers, sincere but systematically mistaken reviewers, qualified dissenters,
Sybils, coordinated credentialed humans, shared-control clusters, honest and
captured roots, and operator capture. Attack intensity, graph sparsity, bridge
coverage, and outcome delay should be varied on a frozen grid.

Primary outcomes are capture radius; false eligibility and exclusion; confirmed
material-defect yield per reviewer-hour; survival of qualified dissent; cross-topic
leakage; control concentration; newcomer wait and supervision burden; appeal and
correction outcomes; measured persona linkage; deterministic replay; fork
localization; and the multiaxis resource vector. Distributions and topic,
newcomer, and minority-view subgroups are mandatory because a global average can
hide capture of a small specialty.

\subsection{Open hypotheses and promotion rules}

\openhypothesis{H1: A3 and A4 reduce capture radius relative to A2 when malicious
influence enters through a bounded frontier.}

\openhypothesis{H2: A4 reduces dense-ring capture relative to A3, with a measurable
coverage cost in sparse domains.}

\openhypothesis{H3: explicit topic/capability scoping and importer-approved bridges
reduce cross-topic false eligibility relative to A2 without unacceptable loss of
legitimate multidisciplinary coverage.}

\openhypothesis{H4: delayed typed outcomes preserve qualified dissent better than
standing derived from raw votes.}

\openhypothesis{H5: forked overlays localize a captured root or policy while
preserving common record identifiers.}

\openhypothesis{H6: brokered per-case personas enforce case-scoped duplicate limits
and private appeals, while declared linkage attacks force privacy-claim downgrade in
some regimes.}

Before outcome inspection, freeze the data-generating process, seed count and its
precision rationale, primary contrasts, multiplicity handling, subgroup harm bounds,
privacy thresholds, failed-run treatment, sensitivity grid, and stop rules. A
hypothesis becomes a result only when the immutable report contains the complete
planned analysis, including null and adverse findings. A synthetic result remains
synthetic.

\begin{evidencegate}[Comparative-result gate]
No A0--A4 code, dataset, seed manifest, or report was present in the inspected v0.1
artifact. This draft therefore contains no comparative result prose.
\end{evidencegate}

\section{Privacy, accountability, and governance limits}

Anonymity is never an unqualified field in the profile. A privacy claim must name
the protected linkage, adversary, operator coalition, auxiliary information,
minimum effective set, attack suite, measurement, and expiry. Timing, rare
attributes, graph topology, prose stylometry, and small pools can defeat a system
whose cryptographic identifiers are otherwise fresh. The correct output under a
failed threshold may be ``confidential under named operators,'' not ``anonymous.''

Accountability creates a second boundary. Duplicate prevention, qualified appeals,
and exceptional opening can require private linkage or credentials. The protocol
therefore specifies narrow purposes and separated roles rather than claiming that
accountability and anonymity coexist automatically. Routine management curiosity,
scientific disagreement, criticism of a senior actor, or an unadjudicated allegation
must not be an opening cause. Corrective events append to history; they do not erase
the original allegation or finding.

Prior anonymous review and reputation systems demonstrate that sophisticated
privacy mechanisms are possible under explicit assumptions
\citep{naessens2006fair,zhai2016anonrep}. \tdrg{} does not reproduce those
mechanisms. A later implementation must either adopt and analyze a suitable
credential/accountability construction or narrow its claim to ordinary
confidentiality. Privacy review must be external to the drafting and implementation
team before a participant pilot.

The graph also cannot regulate actors by itself. Roots, policy authorities,
registrars, evaluators, and appeals bodies are governed institutions. Their funding,
control, conflicts, and replacement procedures belong in the manifest. A local
view can help communities separate from a captured or incompatible domain, but it
cannot make governance costless or guarantee that actors select a trustworthy fork.

\section{Standards interoperability}

\tdrg{} is not a replacement scholarly metadata universe. NISO Z39.106-2023
standardizes terminology for peer-review practices \citep{niso2023peerreview}.
COAR Notify defines interoperable notifications between repositories and review
services using Web protocols and JSON-LD patterns \citep{coar2024notify}.
DocMaps provides machine-readable descriptions of editorial and review processes
\citep{docmaps2021framework}. Crossref accepts metadata for peer-review objects and
relations such as \texttt{isReviewOf} \citep{crossref2026peerreview}. These
infrastructures cover terminology, transport, event description, and registration;
\tdrg{} should profile them rather than duplicate them.

The scientific-record side can reuse general provenance and research-object
standards. PROV-O represents entities, activities, agents, and their provenance
relations for interchange \citep{lebo2013provo}. RO-Crate packages research
artifacts and contextual metadata, including workflows and provenance
\citep{soilandreyes2022rocrate,rocrate2025spec}. These representations do not by
themselves encode the entire trust-domain policy, but they provide established
substrates for artifact identity and lineage.

The normative interoperability requirement is semantic preservation. A projection
must state its source and target versions, extension fields, access restrictions,
and loss boundary. Round-trip fixtures should compare identity scope, claim scope,
review target, challenge and supersession links, and policy meaning. If a public
standard cannot carry a private field, the projection should retain a commitment or
access-class marker rather than disclose the field or silently drop its significance.

\begin{evidencegate}[Interoperability gate]
The v0.1 tooling has no COAR Notify, DocMaps, Crossref, NISO, PROV-O, or RO-Crate
projection and no round-trip fixture. TM-16 is wholly unimplemented.
\end{evidencegate}

\section{Reproducibility and release contract}
\label{sec:repro}

The protocol should be evaluated under the same machine-verifiable record discipline
that it asks of scientific papers. Each result bundle must bind protocol and schema
versions, code revision, environment, configuration, seed stream, input snapshot,
clock, platform, resource vector, command, outputs, test status, and reviewer
sign-off. Public artifacts should have immutable identifiers when permitted;
restricted artifacts should disclose identity commitments, access class, and the
procedure by which an authorized reviewer can inspect them.

\subsection{Multiaxis resources and downselect modes}

Technical burden is represented by R0--R3 for convenience, but the canonical
ten-axis resource vector is authoritative: compute, storage, access, platform,
wall time, human intelligence, agent intelligence, operational support, monetary
or allocative cost, and freshness. Implementations may retain finer subfields such
as memory, CPU versus accelerator time, working versus archival storage, required
roles, and scheduler/facility dependence. Intelligence is therefore a declared
resource rather than an invisible free input, and no additive or fungible scale is
assumed.

R0 denotes documented operation with contemporary portable tools on lightweight
compute and no specialty scheduler. R1 permits a larger workstation or local
accelerator. R2 requires managed HPC, a specialty scheduler, substantial accelerator
use, or comparable support. R3 includes restricted facilities/data, facility-specific
services, or prohibitive working storage. These labels describe access burden, not
scientific quality.

The canonical assessment modes are full execution, bounded recomputation, digest
audit, restricted-platform witnessed attestation, frozen-output inspection, and an
independent evidence path. Witnessed execution and bounded recomputation remain
distinct because custody and inferential scope differ. Every non-full mode must map
included and excluded claims and expose the applicable generator code, input
identities, versions, scientific invariants, sampling rules, attestations, and
residual trust boundaries. A digested product is assessable only within that scope.
No assessment mode is itself an epistemic achievement, and a digest audit must
never be relabeled as full execution or scientific reproduction.

\subsection{Freshness and adverse status}

Every external standard, dependency, credential, data service, and environment has
a freshness date and review cadence. Scheduled reruns should record dependency
drift even when results remain stable. If later execution fails, the original report
remains in history and a visible adverse-status or supersession event records the
failure, triage, affected claims, and recovery status. Silent replacement destroys
the audit trail needed to understand staleness.

The current local 11-test run is R0-like in execution burden, but it is not yet an
archival reproduction package: it lacks an immutable release identifier, independent
rerun, finalized licensing, and a signed or otherwise registered attestation. This
sentence is a limitation, not a request to infer those properties from Git history.

\section{Discussion and stop rules}

The protocol's central wager is that a common record and plural reliance views can
coexist. This structure may reduce the blast radius of a captured policy because an
independent community can recompute from different roots or fork its overlay. It may
also make disagreements legible: two domains can share every review object while
declaring different routing and reliance rules. Those are design affordances, not
measured social outcomes.

The costs are substantial. Fine-grained domains may become sparse; bridge governance
may delay multidisciplinary work; privacy services and appeals require trustworthy
operators; control metadata can itself be sensitive; and aggressive attack defenses
may exclude newcomers or minority specialties. Human and agent intelligence,
governance labor, restricted access, and HPC support must be included in the cost
comparison. A technically elegant graph that cannot be operated or appealed is not
a credible governance mechanism.

Promotion should stop if policy replay is not deterministic from frozen inputs; a
review vote or scientific disagreement can automatically sanction an actor; authority
crosses topics without importer consent; ordinary operators can violate the declared
persona boundary; forks destroy shared record identity; or newcomer/minority false
exclusion exceeds frozen limits. Promotion should also stop if A4 is not materially
better than a simpler arm after accounting for coverage, privacy, and operational
cost. In that case, the useful product is the event vocabulary, conformance report,
and interoperability profile.

The current implementation passes a much smaller test: it makes a local-view rule
inspectable and exercises 11 bounded behaviors. That is enough to prevent the paper
from remaining wholly architectural, but not enough to claim security, anonymity,
fairness, attack resistance, governance quality, or scientific validity. The next
credible artifact is not stronger prose. It is a versioned conformance release with
subtest coverage and an adversarial evaluation frozen before outcome inspection.

\section{Conclusion}

Machine-verifiable scientific review needs both durable records and governed actors.
\tdrg{} proposes to connect them without collapsing all communities, topics, and
times into one reviewer score. It specifies shared typed events, local rooted views,
explicit bridges, forkable overlays, authority boundaries, and strict promotion
rules. Prior work supplies most component mechanisms; the research question is
whether their constrained composition produces a useful and testable governance
boundary.

The bounded v0.1 evidence establishes only that 11 unit tests pass for one private
revision and its synthetic fixtures. Everything else in TM-01--TM-18 remains a
roadmap unless and until a named artifact passes a named gate. This separation is
the protocol's most important present result: machine checks can make claims more
traceable, but they cannot automate the scientific and institutional judgment that
gives review its meaning.

\bibliographystyle{plainnat}
\bibliography{references}
\end{document}
